\begin{resumo}[Abstract]
\begin{otherlanguage*}{english}
\begin{SingleSpacing}
Continuous maintenance of urban infrastructure represents a structural challenge for public administration in medium- and large-sized Brazilian municipalities, where everyday failures --- such as potholes in asphalt paving, non-functional public lighting poles, blockages or leaks in stormwater drainage networks, malfunctioning traffic lights, and irregular solid waste disposal sites --- compromise road safety, environmental health, and collective well-being. In most municipal governments, the information flow between citizens who identify problems in the territory and operational field teams responsible for repairs is fragmented, resulting in duplicate service requests, imprecise geographic location data, lack of standardized visual evidence, and insufficient control over fleet dispatch, workforce allocation, and warehouse material consumption. Given this context, this technical and academic report, corresponding to Stage 1 of the Database course project, aims to describe in detail the problem of collaborative urban infrastructure monitoring and maintenance, substantiate its alignment with the United Nations (UN) 2030 Agenda Sustainable Development Goals (SDGs), and define the functional scope of a relational database system designed to address it. The analysis is grounded in official data from the Brazilian Institute of Geography and Statistics (IBGE), the Institute for Applied Economic Research (IPEA), the National Sanitation Information System (SNIS), and Brazilian federal legislation on urban management and public service user rights. The study demonstrates that the proposed solution aligns directly with SDG 11 (Sustainable Cities and Communities), specifically Targets 11.3, 11.6, and 11.7, and complementarily with SDG 9 (Industry, Innovation and Infrastructure, Target 9.1) and SDG 16 (Peace, Justice and Strong Institutions, Targets 16.6 and 16.7). Finally, the high-level scope of the database is established across twelve business entities organized into three integrated subdomains --- citizen demands geolocated by postal code (CEP), municipal secretariat operational structure, and work order execution linked to the supply chain --- providing the conceptual foundation for the entity-relationship and relational modeling developed in Stage 2.

\vspace{\onelineskip}
\noindent
\textbf{Keywords}: Urban infrastructure. Municipal maintenance. Relational database. Sustainable Development Goals. SDG 11.
\end{SingleSpacing}
\end{otherlanguage*}
\end{resumo}
